\section{Conclusions}
In this thesis, an exocentric vision framework has been deployed for a contact inspection task, through a cooperative dual-UAV \textit{Shared Control} architecture . The framework decouples the overall mission into separate roles: an Inspection UAV designed for safe physical contact, and an Observer UAV tasked with providing an "optimal" and active Point of View (PoV) for a human operator via a bilateral haptic teleoperation interface.
\\To ensure visual tracking and situational awareness, the Observer UAV is governed by a Nonlinear Model Predictive Controller (NMPC) operating in relative cylindrical coordinates. By continuously regulating relative horizontal distance, relative elevation, and camera heading, the optimal solver guarantees that the target vehicle remains centered within the camera Field of View. Furthermore, the explicit penalization of higher-order dynamic derivatives (jerk and snap) within the cost function effectively minimizes aggressive camera motions and operator motion sickness. On the interaction side, the Inspection UAV incorporates a dynamic admittance filter that translates measured or estimated contact forces into compliant reference displacements, preventing mechanical bouncings and ensuring stable physical coupling against vertical structures.
\\A shared teleoperation architecture based on the Novint Falcon haptic device bridges the human operator with the multi-agent system. The device maps pilot inputs to relative viewpoint adjustments while sending tactile safety cues through an Artificial Potential Field (APF). Exponential repulsive forces prevent the operator from violating the minimum inter-UAVs safety distance or driving the target outside the optical aperture.
\\The system has been validated both in high-fidelity Software-In-The-Loop (SITL) simulations and through physical flight experiments. Embedded real-world tests executed onboard a custom octocopter powered by a LattePanda companion computer demonstrated the computational feasibility of the framework: the NMPC solver executed with an average solve time of $2.57\text{ ms}$ at $100\text{ Hz}$, with overruns limited to $0.08\%$. Physical flight tests further validated the Shared Control framework through the bilateral haptic interface: despite a mean one-way wireless communication latency of $45.56\text{ ms}$, the tactile feedback remained completely responsive. When human inputs drove the vehicle towards the minimum standoff distance ($r_{\min} = 1.0\text{ m}$) or the vertical camera FoV limits, the APF triggered repulsive forces up to the $15\text{ N}$ saturation limit, reliably repelling the human operator and restoring a safe observation geometry.
\section{Future Works}
Despite good results already obtained, several directions can be pursued to improve and further advance this research:
\begin{itemize}
    \item \textbf{Simultaneous Dual-UAV Real Flights:} Extending the physical experimental campaign to concurrently deploy both custom UAV platforms in the flight arena, validating the complete cooperative mission pipeline including dynamic contact in hardware.
    \item \textbf{Vision-Based Onboard Target Detection:} Replacing the external motion capture system (OptiTrack) with onboard visual detection algorithms running on the camera payload, closing the perception-action loop entirely onboard.
    \item \textbf{Obstacle Avoidance:} Augmenting the NMPC formulation and the haptic APF with static or dynamic obstacle avoidance constraints to handle complex scenarios with static and moving obstacles.
    \item \textbf{Adaptive Admittance Control:} finding a mathematical definition of Optimal Point Of View, at least for this task.
\end{itemize}
